% !TeX root = ../papers/nc-p3.tex
\section{直線成分の反対縮約による全次数AS判定}
\label{sec:line-component-as-test}
標準有限双対Frobenius窓で各辺の左右のファイバーが空か一点、共通グラフ $X$ が一次元とする。
$X$ の被約成分に直線があれば、定理\ref{thm:point-graph-curve-degree}により全位置で直線である。
重複した台の直線、埋込み点、他の任意の曲線成分を許す。
各位置の文字空間を $V_i$、直線成分 $Z$ の消滅二平面を $U_i(Z)\subset V_i$ と書く。

\begin{theorem}\label{thm:line-component-as-test}
任意の一本の直線成分の四移動 $\sigma$ による軌道 $\mathcal L$ を選ぶ。
位置 $i+2,i+3$ の二文字で四次テンソルを $Z\times Z$ に制限し、
直線の対角方程式 $s_0t_1-s_1t_0$ を除いた二文字関係を $r_i(Z)$ とする。このとき
\[
 0\ne r_i(Z)\in R_i\cap\big(U_i(\sigma Z)\otimes U_{i+1}(Z)\big).
\]
全次数で標準Koszul左右AS正則であるための必要十分条件は、
$i=0,1,2,3$ と全 $Z\in\mathcal L$ について $r_i(Z)$ が階数2を持つことである。
特に各位置で空間を張る既約曲線成分と一本の直線成分が共存すれば、
全ての有限Frobenius候補はAS正則である。
\end{theorem}
\begin{proof}[関係と低次の解消]
$Z$ をどの位置でも同じ $\mathbb P^1$ でパラメータ化する。最後の二文字の各関係の制限は
対角上で消える $(1,1)$ 形式で、その像は一次元である。零なら対応が直線の積を含み点ファイバーに反する。
$w$ の六対六対合から $r_i(Z)\ne0$ である。第三文字を $s$ で評価し、第二文字も $s$ で評価すれば
隣接関係により零となる。$t\ne s$ を選べば $r_i(Z)$ の第二文字は $U_{i+1}(Z)$ に入る。
同様に第四文字を $t$、最初を $\sigma t$ と評価すれば、最後から最初への関係で第一の包含も得る。

$D(Z)$ を直線の切断環とする。全位置で $\mathcal O(1)$ の二切断が得られるので、
$A\to D(Z)$ は全射で、$D(Z)$ は二変数多項式環の反復である。
定理\ref{thm:elliptic-quartic-as-test}の束複体を $\mathbb P^1$ に制限する。
内部次数 $n$ で項 $j$ の束は $\mathcal O(n-j)$ の $1,4,6,4,1$ 個の直和である。
点評価の階数は $1,3,3,1$ で局所分裂完全である。$n\ge3$ では全項の $H^1$ が零、
$n=2$ では項4に $H^1(\mathcal O(-2))=k$ が残り、二行の超コホモロジーの微分で
切断列の項2にちょうど一つのホモロジーを与える。$n=1$ では切断列自体が四文字から二切断への
全射なので項1に二つ、$n=0$ は項0のみである。反対側も同じである。
単純の極小解消の初めの二微分は候補と同一、三次の重なりも第三項に入るので、
実際の第一Torはこの計算と一致し、第二Torは計算した一元の商となる。
従って $D(Z)$ の始点 $i$ での右イデアルは $U_i(Z)A$ で、
その関係は次数2で高々一つ、左イデアルも各終点で同様である。
既に非零な $r_i(\sigma^{-1} Z)$ がその関係を与えるので必ず一つである。

\smallskip
\noindent\textit{必要性。}
標準ASなら候補は解消であり、直線商の左右極小解消は次数 $0,1,2$、個数 $1,2,1$ で終わる。
もし最後の二係数が比例するなら、右解消を双対化した第二Extの最小生成元に対する
左側の一次消滅空間は高々一次元である。一方、軌道長の四倍の周期に折り返せば、
定理\ref{thm:elliptic-quartic-as-test}の局所双対計算がそのまま適用され、
この第二Extは標準四移動を除いて $R^2\Gamma_A(D(Z))'$ に一致する。
その作用はひねられた直線商を通るので一次消滅空間は少なくとも二次元である。
実際ここで使う直線商の左右解消は有限自由であり、$A/A_{\ge N}$ との正Torは
内部次数 $[N,N+1]$ に限られる。従って同定の証明で用いた入射加群の非輪状性も同様に成立する。
これで階数1は排除される。

\smallskip
\noindent\textit{十分性。}
全整数で $Z_{i+1}=\sigma^{-1}Z_i$ とし、二文字平面 $U_i(Z_i)$ と
非退化関係 $r_i(Z_{i+1})$ で抽象的な部分表示 $P$ を作る。各位置で順に基底を取れば
関係が交換子となり $H_P=(1-t)^{-2}$ である。
$U_i(Z_i)A$ による右商は $D(Z_i)$ の始点部分、$A U_{j-1}(Z_{j-1})$ による左商は
$D(Z_{j-1})$ の終点部分で、いずれも長さ $n$ で次元 $n+1$ である。
各商の基底を持ち上げて $P$ の語を前または後ろに掛ければ $A$ を張る。
従って各成分の次元は $(1-t)^{-4}$ の係数以下である。一方、右表示の唯一の関係で始めた列のEuler式は、
固定終点で $a_n-2a_{n-1}+a_{n-2}=n+1+e_n$（$e_n\ge0$ は末項の核次元）である。
$(1-t)^{-2}$ を掛ければ上界で等号、全 $e_n=0$ である。左右の張る集合は基底で、
$A$ は $P$ 上左右自由、始点 $i$ の直線商加群は $P$ 単純の誘導となる。
任意の所定の $i,Z$ に対し $Z_i=Z$ としてこの構成を行えるから、同じ $D(Z)$ の全始点で
$\underline{\operatorname{Ext}}^j_A(e_iD(Z),A)=0$ ($j<2$) を得る。終点側も同じである。
直線商を通る環変更で左右の単純から $A$ への最初の二Extも消える。
生成と二次表示で候補複体の始めは完全であり、上の消滅とFrobenius双対で末尾も完全、
中央は既に求めたHilbert関数で完全である。従って全次数のKoszul性と左右AS性が従う。

最後に $r_i(Z)=a\otimes b$ ならグラフの各既約成分で最初の線型形式 $a$ または二番目の
線型形式 $b$ が恒等零となる。両位置で空間を張る同じ既約成分があれば不可能である。
\end{proof}

\begin{corollary}\label{cor:line-component-atlas}
直線を含む一次元点グラフの全AS窓は、次の有限パラメータと同型判定で被覆される。
最初の三位置の指定直線を同じ座標直線にして、その上の隣接写像を恒等とする。
二つの辺ごとに、直線上の切断乗法核である十三次元空間の六平面を選ぶ。
二窓からの有限Frobenius復元、四辺の点ファイバー、一次元の像と上の軌道の階数検査を課す。
指定を忘れた同型は、第二の窓の有限個の直線成分のどれかを指定して、
三位置の線型変換によって二つの関係を同時に運べることと同値である。
これらの変換は当該直線を対応させ、そこで同じ $\mathrm{PGL}_2$ を誘導するものに限れば十分である。
同じ比較を一移動・二移動に用いれば最小周期を得る。
\end{corollary}
\begin{proof}
全窓で一つの初期直線は点グラフによって一意に運ばれる。成分ごとの次数一定性で
それは各位置の直線で、全次数は高々7だから軌道探索は有限である。
逆に最初の二辺に指定直線のグラフが入り、復元・ファイバー・次元検査を通れば、
共通像の既約成分としてその直線が入る。二窓の延長は系\ref{cor:two-window-extension}による。
座標変更の存在はこの三つの放物部分群と共通の二次行列を用いた有限な多項式条件で判定できる。
\end{proof}
